\section{Introducción}

Este documento es una guía de estudio sobre ROS~2 (Robot Operating System 2) orientada a explicar el \textit{por qué} de cada decisión de diseño: por qué existe ROS~2 si ya existía ROS~1, por qué se eligió DDS como middleware, por qué el ecosistema está atado a versiones específicas de Ubuntu, y por qué distintas herramientas de simulación tienen distintos niveles de integración.

El documento está organizado de lo general a lo particular. Las primeras secciones cubren el contexto histórico: cómo nació ROS, qué limitaciones llevaron al rediseño, y cómo ese rediseño tomó forma en ROS~2. Las secciones siguientes bajan al nivel técnico del ecosistema: el middleware DDS, el esquema de distribuciones, los simuladores disponibles, y la arquitectura fundamental de nodos. El conjunto forma el mapa conceptual necesario para trabajar con ROS~2 de forma informada, no solo mecánica.

La motivación práctica de este documento es el trabajo de experimentación en control de formación multi-robot sobre ROS~2 Humble con Gazebo Classic 11. Comprender el ecosistema en el que se trabaja permite tomar decisiones de diseño fundamentadas y depurar problemas en el nivel correcto.
